\begin{lemma}[Idempotence in terms of the diagonal matrices]%
	\label{lem:mpdo_bnt_idempotent_coefficient_form_of_chi}
	\lean{MPOTensor.BNTAlgebraClause.idempotent_coefficient_form_ofChi}
	\leanok
	\uses{def:mpdo_bnt_algebra_clause}
	Suppose that the BNT algebra conditions hold. The vector $m$ is then an
	idempotent for the multiplication induced by the length-one
	coefficients $\tr(\chi_{\alpha,\beta,\gamma})$: for every $\gamma$,
	\begin{align}
		m_\gamma
		 & =\sum_{\alpha,\beta}
		\tr(\chi_{\alpha,\beta,\gamma})
		m_\alpha m_\beta.
		\notag
	\end{align}
\end{lemma}

\begin{proof}\leanok
	The trace-power representation at length one gives
	$c^{(1)}_{\alpha,\beta,\gamma}
		=\tr(\chi_{\alpha,\beta,\gamma})$.
	Substitution into the idempotent law
	$m_\gamma
		=\sum_{\alpha,\beta}
		c^{(1)}_{\alpha,\beta,\gamma}m_\alpha m_\beta$
	gives the stated identity.
\end{proof}

\begin{definition}[Multiplicity spectra from eventual equality of power sums]%
	\label{def:mpdo_bnt_eventual_power_sums_multiplicity_comparison}
	\lean{MPOTensor.BNTMultiplicitySpectrumComparison.%
		ofEventuallyEqualPowerSums}
	\leanok
	\uses{def:mpdo_bnt_multiplicity_spectrum_comparison, lem:eventual_power_sum_equality_all_exponents, thm:bounded_power_sum_multiset}
	Let $\mu_\alpha$ and $\nu_\gamma$ be the diagonal multiplicity matrices
	in the one-site and two-site vertical decompositions, respectively.
	Suppose that their diagonal entries and those of
	$\chi_{\alpha,\beta,\gamma}$ are strictly positive, and that
	$\sum_i\mu_{\alpha,i}=m_\alpha$ for every $\alpha$.
	If, for every $\gamma$ and all sufficiently large $L$,
	\begin{align}
		\sum_r \nu_{\gamma,r}^{L}
		 & =
		\sum_{\alpha,\beta,i,j,k}
		\bigl(\mu_{\alpha,i}\mu_{\beta,j}
		\chi_{\alpha,\beta,\gamma,k}\bigr)^L,
		\notag
	\end{align}
	then the multiset of entries $\nu_{\gamma,r}$ equals the multiset of
	products $\mu_{\alpha,i}\mu_{\beta,j}\chi_{\alpha,\beta,\gamma,k}$.
	These multiplicity matrices, together with this multiset equality and
	the one-site trace identity, give a comparison of the one-site and
	two-site multiplicity spectra.

	Positivity makes all entries nonzero. Eventual equality of two finite
	sums of geometric sequences gives, for every $L\geq 1$,
	\begin{align}
		\sum_r \nu_{\gamma,r}^{L}
		 & =
		\sum_{\alpha,\beta,i,j,k}
		\bigl(\mu_{\alpha,i}\mu_{\beta,j}
		\chi_{\alpha,\beta,\gamma,k}\bigr)^L.
		\notag
	\end{align}
	The finite power-sum identity determines both the cardinality and the
	multiset of entries.
\end{definition}

\begin{theorem}[Multiplicity trace normalization]%
	\label{thm:mpdo_bnt_multiplicity_trace_normalization}
	\lean{MPOTensor.BNTMultiplicitySpectrumComparison.sum_twoEntry_eq_sum_products}
	\lean{MPOTensor.BNTMultiplicitySpectrumComparison.sum_products_eq}
	\lean{MPOTensor.BNTMultiplicitySpectrumComparison.twoTrace_eq_traceScalar}
	\lean{MPOTensor.BNTAlgebraClause.twoMultiplicityTrace_eq_traceScalar}
	\leanok
	\uses{def:mpdo_bnt_algebra_clause, def:mpdo_bnt_multiplicity_spectrum_comparison, lem:mpdo_bnt_idempotent_coefficient_form_of_chi}
	Suppose that the BNT algebra conditions hold and that the one-site and
	two-site vertical multiplicity spectra satisfy the multiset comparison
	just defined. Then every two-site multiplicity matrix has trace
	$m_\gamma$:
	\begin{align}
		\tr(\nu_\gamma)
		 & =\sum_{\alpha,\beta,i,j,k}
		\mu_{\alpha,i}\mu_{\beta,j}
		\chi_{\alpha,\beta,\gamma,k}
		=\sum_{\alpha,\beta}
		c^{(1)}_{\alpha,\beta,\gamma}m_\alpha m_\beta
		=m_\gamma.
		\label{eq:rfp_trace_normalization}
	\end{align}
	This is the normalization used for the preparation matrix
	$\nu_\gamma/m_\gamma$ in the algebra-to-RFP implication of
	\cite[Appendix~C.4, lines~2058--2064]{Cirac2016MPDO_arXiv}.
\end{theorem}
\begin{proof}\leanok
	Sum the multiset equality of the two spectra. For each
	$\alpha,\beta,\gamma$, the resulting product sum satisfies
	\begin{align}
		\sum_{i,j,k}
		\mu_{\alpha,i}\mu_{\beta,j}\chi_{\alpha,\beta,\gamma,k}
		 & =
		\left(\sum_i\mu_{\alpha,i}\right)
		\left(\sum_j\mu_{\beta,j}\right)
		\left(\sum_k\chi_{\alpha,\beta,\gamma,k}\right).
		\notag
	\end{align}
	The identities
	\begin{align}
		\sum_i\mu_{\alpha,i}
		 & =m_\alpha,
		\notag                            \\
		\sum_j\mu_{\beta,j}
		 & =m_\beta,
		\notag                            \\
		\sum_k\chi_{\alpha,\beta,\gamma,k}
		 & =c^{(1)}_{\alpha,\beta,\gamma}
		\notag
	\end{align}
	therefore give
	\begin{align}
		\tr(\nu_\gamma)
		 & =\sum_{\alpha,\beta}
		m_\alpha m_\beta c^{(1)}_{\alpha,\beta,\gamma}
		=m_\gamma,
		\notag
	\end{align}
	where the final equality is the length-one idempotent law.
\end{proof}

\begin{definition}[BNT algebra with a blocked-basis comparison]%
	\label{def:mpdo_bnt_label_theorem_data}
	\lean{MPOTensor.BNTLabelTheoremData}
	\lean{MPOTensor.BNTLabelTheoremData.same_length_product_form}
	\lean{MPOTensor.BNTLabelTheoremData.idempotent_coefficient_form}
	\lean{MPOTensor.BNTLabelTheoremData.positiveChi}
	\lean{MPOTensor.BNTLabelTheoremData.same_length_product_form_ofChi}
	\lean{MPOTensor.BNTLabelTheoremData.idempotent_coefficient_form_ofChi}
	\lean{MPOTensor.BNTLabelTheoremData.same_length_product_eq_sum_ofChi}
	\lean{MPOTensor.BNTLabelTheoremData.idempotent_eq_sum_ofChi}
	\lean{MPOTensor.BNTLabelTheoremData.positive_chi_pos_entries}
	\lean{MPOTensor.BNTLabelTheoremData.positive_chi_trace_power}
	\lean{MPOTensor.BNTLabelTheoremData.labelAssignment}
	\lean{MPOTensor.BNTLabelTheoremData.sourceLabel}
	\lean{MPOTensor.BNTLabelTheoremData.targetLabel}
	\lean{MPOTensor.BNTLabelTheoremData.positiveBlockedChi}
	\lean{MPOTensor.BNTLabelTheoremData.ofChi}
	\lean{MPOTensor.BNTLabelTheoremData.same_length_product_eq_sum}
	\lean{MPOTensor.BNTLabelTheoremData.idempotent_eq_sum}
	\lean{MPOTensor.BNTLabelTheoremData.blocked_coeff_eq}
	\lean{MPOTensor.BNTLabelTheoremData.blockedComparison_ofChi}
	\lean{MPOTensor.BNTLabelTheoremData.blocked_coeff_eq_ofChi}
	\leanok
	\uses{def:mpdo_bnt_label_coefficient_family, def:mpdo_bnt_label_operator_family, def:mpdo_bnt_label_trace_scalar_family, def:mpdo_bnt_algebra_clause, def:mpdo_bnt_blocked_basis_label_assignment, def:mpdo_bnt_blocked_basis_coefficient_comparison}
	Consider operators $O_L(M_\alpha)$, structure coefficients
	$c^{(L)}_{\alpha,\beta,\gamma}$, and scalars $m_\alpha$ indexed by a
	finite set of BNT labels. In the vertical MPDO decomposition,
	$m_\alpha=\tr(\mu_\alpha)$.
	Require the BNT algebra conditions and a comparison with chosen
	blocked bases.
	The algebra conditions require positive diagonal matrices
	$\chi_{\alpha,\beta,\gamma}$ such that
	$c^{(L)}_{\alpha,\beta,\gamma}
		=\tr(\chi_{\alpha,\beta,\gamma}^{\,L})$ for every $L>0$,
	and, at every such length, the product law
	\begin{align}
		O_L(M_\alpha)O_L(M_\beta)
		 & =\sum_\gamma
		c^{(L)}_{\alpha,\beta,\gamma}O_L(M_\gamma).
		\notag
	\end{align}
	The vector $m$ is an idempotent for the multiplication induced by
	$c^{(1)}$, meaning that, for every $\gamma$,
	\begin{align}
		m_\gamma
		 & =\sum_{\alpha,\beta}
		c^{(1)}_{\alpha,\beta,\gamma}m_\alpha m_\beta.
		\notag
	\end{align}
	The blocked-basis comparison is an additional condition.
	For every positive blocking length $n$, let $\sigma_n$ assign BNT
	labels to the basis indices at length $n$, and let $\tau_n$ assign
	BNT labels to the basis indices at length $2n$.
	For indices $i,j$ at length $n$ and $k$ at length $2n$, require
	\begin{align}
		c^{(n)}_{i,j,k}
		 & =c^{(n)}_{\sigma_n(i),\sigma_n(j),\tau_n(k)}.
		\notag
	\end{align}
	Here the left-hand side denotes the structure coefficient in the chosen
	blocked bases, while the right-hand side uses BNT labels.
	The same maps assign the diagonal matrix
	$\chi_{\sigma_n(i),\sigma_n(j),\tau_n(k)}$ to each triple of
	blocked-basis indices.

	One may start with positive diagonal matrices $\chi_{\alpha,\beta,\gamma}$
	and define the coefficients by their traces of powers. In that case,
	the product law, idempotent law, and blocked-basis comparison are
	required for those coefficients. Conversely, replacing coefficients
	satisfying the trace-power identity by
	$\tr(\chi_{\alpha,\beta,\gamma}^{\,L})$ preserves these conditions at
	every positive length.
\end{definition}

\begin{theorem}[BNT structure coefficients and product formulas]%
	\label{thm:mpdo_bnt_label_theorem_data_source_formulas}
	\lean{MPOTensor.BNTLabelTheoremData.coeff_eq_trace_pow}
	\lean{MPOTensor.BNTLabelTheoremData.coeff_eq_ofChi_coeff}
	\lean{MPOTensor.BNTLabelTheoremData.chi_entry_pos}
	\lean{MPOTensor.BNTLabelTheoremData.same_length_product_eq_sum_chi_trace_pow}
	\lean{MPOTensor.BNTLabelTheoremData.idempotent_eq_sum_chi_trace}
	\leanok
	\uses{def:mpdo_bnt_label_theorem_data}
	Suppose that the BNT algebra conditions and blocked-basis comparison
	hold. Let $\chi_{\alpha,\beta,\gamma}$ be the associated positive
	diagonal matrices. For every $L>0$,
	\begin{align}
		c^{(L)}_{\alpha,\beta,\gamma}
		 & =c^{\chi,(L)}_{\alpha,\beta,\gamma}
		=\tr\!\left(\chi_{\alpha,\beta,\gamma}^{\,L}\right),
		\notag
	\end{align}
	where $c^{\chi,(L)}_{\alpha,\beta,\gamma}$ denotes the trace of the
	$L$th power of $\chi_{\alpha,\beta,\gamma}$. Every diagonal entry of
	every $\chi_{\alpha,\beta,\gamma}$ is strictly positive.
	Consequently, for every $L>0$, the product identity reads
	\begin{align}
		O_L(M_\alpha)O_L(M_\beta)
		 & =\sum_\gamma
		\tr\!\left(\chi_{\alpha,\beta,\gamma}^{\,L}\right)
		O_L(M_\gamma),
		\notag
	\end{align}
	and the idempotent identity reads
	\begin{align}
		m_\gamma
		 & =\sum_{\alpha,\beta}
		\tr(\chi_{\alpha,\beta,\gamma})
		m_\alpha m_\beta.
		\notag
	\end{align}
\end{theorem}
\begin{proof}\leanok
	The assumed positive diagonal matrices satisfy
	$c^{(L)}_{\alpha,\beta,\gamma}
		=\tr(\chi_{\alpha,\beta,\gamma}^{\,L})$ for $L>0$.
	Substitution into the product law, and into the idempotent law at
	length one, gives the stated formulas.
\end{proof}

\begin{theorem}[Positive diagonal matrices for the blocked coefficients]%
	\label{thm:mpdo_bnt_label_theorem_data_to_positive_blocked_chi_trace_power_form}
	\lean{MPOTensor.BNTLabelTheoremData.toPositiveBlockedStructureChiTracePowerForm}
	\lean{MPOTensor.BNTLabelTheoremData.positiveBlockedChi_toDiagonal_of_pos}
	\lean{MPOTensor.BNTLabelTheoremData.positiveBlockedChi_tracePowerCoeff_of_pos}
	\lean{MPOTensor.BNTLabelTheoremData.positiveBlockedChi_dim_of_pos}
	\lean{MPOTensor.BNTLabelTheoremData.positiveBlockedChi_trace_matrix_pow_of_pos}
	\lean{MPOTensor.BNTLabelTheoremData.positiveBlockedChi_tracePower}
	\lean{MPOTensor.BNTLabelTheoremData.positiveBlockedChi_posEntries}
	\lean{MPOTensor.BNTLabelTheoremData.positiveBlockedChi_entry_pos}
	\leanok
	\uses{def:mpdo_bnt_label_theorem_data, thm:mpdo_bnt_blocked_basis_to_positive_blocked_chi_trace_power_form}
	Suppose that the BNT algebra conditions and blocked-basis comparison
	hold. Positive diagonal matrices representing the blocked coefficients
	as traces of powers are obtained by pulling back the
	length-independent matrices $\chi_{\alpha,\beta,\gamma}$ along the
	comparison maps. At every positive blocking length $n$, define
	$\chi_{n,i,j,k}
		=\chi_{\sigma_n(i),\sigma_n(j),\tau_n(k)}$.
	Thus the two matrices have equal size, diagonal entries, and traces of all
	powers.
\end{theorem}
\begin{proof}\leanok
	Apply the construction of positive diagonal matrices for blocked
	coefficients to the comparison maps and the positive diagonal matrices
	$\chi_{\alpha,\beta,\gamma}$ representing the BNT structure coefficients.
\end{proof}

\begin{corollary}[Blocked coefficient trace-power formulas]%
	\label{cor:mpdo_bnt_label_theorem_data_blocked_coeff_eq_positive_blocked_chi_trace_pow}
	\lean{MPOTensor.BNTLabelTheoremData.blocked_coeff_eq_trace_pow}
	\lean{MPOTensor.BNTLabelTheoremData.blocked_coeff_eq_positiveBlockedChi_trace_pow}
	\leanok
	\uses{def:mpdo_bnt_label_theorem_data, thm:mpdo_bnt_label_theorem_data_to_positive_blocked_chi_trace_power_form, lem:mpdo_blocked_structure_chi_eq_trace_matrix_pow}
	Under the BNT algebra conditions and blocked-basis comparison, for every
	positive blocking length $n$, the blocked-basis coefficients are
	traces of powers of the matrices $\chi_{\alpha,\beta,\gamma}$ selected
	by the source and target label maps. Equivalently, they are traces of
	powers of the pulled-back matrices $\chi_{n,i,j,k}$:
	\begin{align}
		c^{(n)}_{i,j,k}
		 & =
		\tr\!\left(\chi_{\sigma_n(i),\sigma_n(j),\tau_n(k)}^{\,n}\right)
		=\tr(\chi_{n,i,j,k}^{\,n}).
		\notag
	\end{align}
\end{corollary}
\begin{proof}\leanok
	The blocked-basis comparison and the trace-power representation identify
	the coefficient with the trace of the $n$th power of the matrix
	selected by the source and target maps. By definition, the pulled-back
	matrix $\chi_{n,i,j,k}$ is this same matrix.
\end{proof}
